Die Atommassen in $\mathrm{u}$ stimmen zahlenmässig mit den molaren Massen in $\mathrm{g/mol}$ überein. Nach der Isotopentabelle (FoTaBe S.~199) ist $M_{235} = \MUaO$ und $M_{238} = \MUbO$. Das Verhältnis der Atommassen ist gleich dem Verhältnis der molaren Massen, also
\begin{align}
 \frac{\Delta m}{m_{235}} &= \dmrF \\
   &= \frac{\MUbO-\MUaO}{\MUaO} \\
   &= \dmr \approx \result{\dmrP{0}{2}}
\end{align}
U-238 ist also um rund $\dmrP{0}{2}$ schwerer als U-235. Bezogen auf U-238 ergibt sich auf zwei Stellen derselbe Wert; schon die Nukleonenzahlen (drei Nukleonen mehr bei 235) liefern diese Abschätzung.
